% s-007-p007-m002; original PDF 7, 8, 9
\PrefaceSection{s-007-p007-m002}{3}
Nun aber, was deutet der Name? Sant’ Ilario, der heilige Heiterling oder frohe Heilige: wer mit dem Namen dieses Patrons ein Buch tauft, sollte er wohl die Hand nach zwei Kronen der Vollkommenheit ausstrecken, die bisher getrennt vergeben wurden? Bei der Heiligkeit denkt man an die Wüste, an Gliederstarre und Hungerkuren, an vierzigtägiges oder vierzigjähriges Stillstehen auf einer Säule von Monomanie, oder, wenn man ganz modern sein will, ans Priesterthum der Erkenntniss, an die siegreich wehende Fahne der Wahrheit, unter der der Fahnenträger verwundet zusammenbricht. Die Heiterkeit aber, die kennen wir als überflinke Beweglichkeit, als Feuerwerk oder Schellengeklirr von Narrenkappen, als klapperndes klimperndes klingelndes Scherzo in der Symphonie des Lebens. Wie, wenn einmal der Versuch gemacht würde, heiter zu sein ohne Gezappel und nachdenklich ohne Steifigkeit? wenn man eine unwahrscheinliche Mischung von Nüchternheit und Rausch, von Gauklertanz und Fakirschlaf ausfände und festhielte? Aber das klingt wie ein Programm, und dies Buch hat keines; es füllt keine Lücke aus, kommt keinem tiefgefühlten Bedürfnisse entgegen, es ist stolz darauf, überflüssig zu sein und zu überflüssigen, müssiggehenden, von Arbeit und Amt noch nicht zerriebenen und verbrauchten Menschen zu reden. Dies Buch will alles vermeiden, was wie Zweck, Formel, Willenskrampf aussieht, jede Art Absichtlichkeit im Liegen, Stehen, Springen, Tanzen; giebt es nicht das beste Portrait, wenn man gar kein bestimmtes Gesicht macht? Die hohen Augenbraunen des Erkennenden, die Pose des Spötters, die berufsmässigen Falten im Priesterantlitz oder Galgengesicht — fort mit alledem! Denke, aber ohne ans Denken zu denken! vergiss, und vergiss auch, dass du vergessen wolltest! sei frei, mit der Freiheit, dich ihrer begeben zu können! so ungefähr wird das Buch zu seinem Autor, oder der Autor zum Buche gesprochen haben, ehe Eins vom Anderen Abschied nahm. Dieser Abschied des Neugeborenen von seiner Mutter trug sich in einer jener Landschaften mittlerer Höhe zu, deren Stil ich mit meinem Buche getroffen haben möchte: ein Stil zwischen Pathos und Idyll, eine Landschaft gleich fern von der Niederung wie von fratzenhaft emporgethürmten Bergspitzen. Es war bei der Kirche Sant’ Ilario, die unweit des genuesischen Nervi unter Ölbäumen und Cypressen glänzt und den blauschimmernden Golf von Genua bis hin zum Vorgebirge von Portofino überschaut. Der Namenstag des Heiligen wurde durch eine feierlich-närrische Eselsegnung begangen, wobei unser Aller Wunsch, dass das Eselgeschlecht nicht aussterben möge, in verständlicher Symbolik zur Sprache kam. Der Autor versäumte nicht, sich selbst auf die Seite der Eselinnen zu stellen und für sein Eselsfüllen von Buch den Segen des frohgelaunten Heiligen zu erflehen.
\par

\Sperr{S. Ilario} bei Genua, am Tage S. Ilario 1897.
\par

